\section*{الفصل \ref{ch:g12:trigo} --- الدوال المثلثية}

\begin{solution}{exo:g12:trigo:1}
$\cos\frac{2\pi}{3} = \cos\left(\pi - \frac\pi3\right) = -\cos\frac\pi3
= -\frac12$.

$\sin\frac{5\pi}{6} = \sin\left(\pi - \frac\pi6\right) = \sin\frac\pi6
= \frac12$.

$\cos\frac{7\pi}{4} = \cos\left(2\pi - \frac{\pi}{4}\right)
= \cos\left(-\frac\pi4\right) = \frac{\sqrt2}{2}$.

$\sin\left(-\frac\pi3\right) = -\sin\frac\pi3 = -\frac{\sqrt3}{2}$.
\end{solution}

\begin{solution}{exo:g12:trigo:2}
$f'(x) = 2\sin x\cos x = \sin 2x$ (بقاعدة السلسلة).

$g'(x) = -3\sin(3x) + \sin x + x\cos x$ (بقاعدتي السلسلة والجداء).

بقاعدة خارج القسمة:
\[
h'(x) = \frac{\cos x\,(2 + \cos x) - \sin x\,(-\sin x)}{(2+\cos x)^2}
= \frac{2\cos x + \cos^2 x + \sin^2 x}{(2+\cos x)^2}
= \frac{2\cos x + 1}{(2+\cos x)^2}.
\]
\end{solution}

\begin{solution}{exo:g12:trigo:3}
\emph{(أ)} حل خاص \omterm{def:g10:algebra:equation}{للمعادلة} $\sin x = \frac{\sqrt3}{2}$ هو
$\frac\pi3$. والحلول العامة: $x = \frac\pi3 + 2k\pi$ أو
$x = \pi - \frac\pi3 + 2k\pi = \frac{2\pi}{3} + 2k\pi$. وفي
$\intoc{-\pi}{\pi}$: $x \in \left\{\frac\pi3, \frac{2\pi}{3}\right\}$.

\emph{(ب)} $\cos 2x = \frac{\sqrt2}{2} = \cos\frac\pi4$ يعطي
$2x = \pm\frac\pi4 + 2k\pi$، إذن $x = \pm\frac\pi8 + k\pi$. وفي
$\intoc{-\pi}{\pi}$:
$x \in \left\{-\frac{7\pi}{8}, -\frac\pi8, \frac\pi8, \frac{7\pi}{8}\right\}$.
\end{solution}

\begin{solution}{exo:g12:trigo:4}
مع $a = \frac\pi3$ و $b = \frac\pi4$:
\[
\cos\frac{\pi}{12} = \cos a\cos b + \sin a \sin b
= \frac12\cdot\frac{\sqrt2}{2} + \frac{\sqrt3}{2}\cdot\frac{\sqrt2}{2}
= \frac{\sqrt2 + \sqrt6}{4},
\]
\[
\sin\frac{\pi}{12} = \sin a\cos b - \cos a \sin b
= \frac{\sqrt3}{2}\cdot\frac{\sqrt2}{2} - \frac12\cdot\frac{\sqrt2}{2}
= \frac{\sqrt6 - \sqrt2}{4}.
\]
\end{solution}

\begin{solution}{exo:g12:trigo:5}
ضع $s = \sin x$: فيكون $2s^2 - s - 1 = (2s + 1)(s - 1) \geq 0$ إذا وفقط إذا كان
$s \leq -\frac12$ أو $s = 1$.

و $\sin x = 1$ يعطي $x = \frac\pi2$. وأما $\sin x \leq -\frac12$: فعلى دائرة
الوحدة، هو القوس الواقع تحت المستقيم الأفقي $y = -\frac12$، وهو في
$\intcc{0}{2\pi}$ يساوي $x \in \intcc{\frac{7\pi}{6}}{\frac{11\pi}{6}}$.
ومنه
\[
S = \left\{\frac{\pi}{2}\right\} \cup
\intcc{\frac{7\pi}{6}}{\frac{11\pi}{6}} .
\]
\end{solution}

\begin{solution}{exo:g12:trigo:6}
$\dfrac{\sin 3x}{x} = 3\,\dfrac{\sin 3x}{3x} \to 3 \cdot 1 = 3$
(بالتركيب مع $u = 3x \to 0$).

وباستعمال $1 - \cos x = \dfrac{\sin^2 x}{1 + \cos x}$:
\[
\frac{1 - \cos x}{x^2}
= \left(\frac{\sin x}{x}\right)^{\!2} \frac{1}{1 + \cos x}
\xrightarrow[x\to0]{} 1 \cdot \frac12 = \frac12 .
\]

وبالتعويض $u = \frac1x \to 0^+$:
$x \sin\frac1x = \frac{\sin u}{u} \to 1$.
\end{solution}

\begin{solution}{exo:g12:trigo:7}
$f'(x) = 1 - \sin x \geq 0$ من أجل كل $x$، مع المساواة بالضبط عندما يكون
$\sin x = 1$، \emph{أي} $x = \frac\pi2 + 2k\pi$. وعلى $\intcc{0}{2\pi}$، تتزايد $f$
من $f(0) = 1$ إلى $f(2\pi) = 2\pi + 1$ (وصفر $f'$ عند
$\frac{\pi}{2}$ معزول، فالتزايد تام حتى)؛ \omterm{def:g10:functions:extrema}{والقيمة الصغرى} $1$ عند
$0$، والعظمى $2\pi + 1$ عند $2\pi$. وعلى $\R$، يكون $f' \geq 0$ في كل مكان مع أصفار
معزولة فقط، إذن $f$ \omterm{def:g12:seq:monotonic}{متزايدة} (تمامًا) حسب
\cref{thm:g12:deriv:variations}، مع أن $f'$ تنعدم عند كل
$\frac\pi2 + 2k\pi$.
\end{solution}

\begin{solution}{exo:g12:trigo:8}
\emph{1.} ليكن $R = \sqrt{A^2 + B^2} > 0$. فالنقطة
$\left(\frac AR, \frac BR\right)$ تقع على دائرة الوحدة، إذن يوجد
$\varphi$ حيث $\cos\varphi = \frac AR$ و $\sin\varphi = \frac BR$. عندئذٍ
\[
R\cos(x - \varphi) = R\bigl(\cos x\cos\varphi + \sin x \sin\varphi\bigr)
= A\cos x + B\sin x = f(x).
\]

\emph{2.} وبما أن مجموعة قيم $\cos$ هي $\intcc{-1}{1}$، فللدالة $f$ \omterm{def:g10:functions:extrema}{قيمة عظمى} $R$ و
\omterm{def:g10:functions:extrema}{قيمة صغرى} $-R$. ومن أجل $\cos x + \sin x$: يكون $A = B = 1$ و $R = \sqrt2$ و
$\varphi = \frac\pi4$، فتصير \omterm{def:g10:algebra:equation}{المعادلة}
$\sqrt2\cos\left(x - \frac\pi4\right) = 1$، \emph{أي}
$\cos\left(x - \frac\pi4\right) = \frac{\sqrt2}{2}$، مما يعطي
$x - \frac\pi4 = \pm\frac\pi4 + 2k\pi$: وفي $\intoc{-\pi}{\pi}$،
$x \in \left\{0, \frac\pi2\right\}$.
\end{solution}

\begin{solution}{exo:g12:trigo:9}
$(\sin)'' = (\cos)' = -\sin$ و $(\cos)'' = (-\sin)' = -\cos$: فكلتاهما تحقق
$y'' = -y$.

ولتكن $g = y - a\cos - b\sin$. عندئذٍ $g'' = -g$ (بخطية الاشتقاق)، و
$g(0) = y(0) - a = 0$ و $g'(0) = y'(0) - b = 0$. ضع
$E = g^2 + (g')^2$. عندئذٍ
\[
E' = 2g g' + 2g' g'' = 2g g' - 2g' g = 0,
\]
إذن $E$ ثابتة تساوي $E(0) = 0$. ومجموع مربعين ينعدم
تطابقيًا يفرض $g = 0$، ومنه $y = a\cos + b\sin$.
\end{solution}

\begin{solution}{pb:g12:trigo:1}
\textbf{1.} $0.99833\ldots$ و $0.99998\ldots$: أي زحفًا نحو
$1$.

\textbf{2.} المثلث $OAM$: القاعدة $1$ والارتفاع $\sin x$: فمساحته
$\frac{\sin x}{2}$. والقطاع $OAM$: وهو جزء
$\frac{x}{2\pi}$ من قرص الوحدة: فمساحته $\frac x2$. والمثلث
$OAT$: القاعدة $1$ والارتفاع $\tan x$: فمساحته $\frac{\tan x}{2}$.

\textbf{3.} المثلث يجلس داخل القطاع الجالس داخل المثلث
الكبير: $\frac{\sin x}{2} < \frac x2 < \frac{\tan x}{2}$،
أي $\sin x < x < \tan x$. ومن $\sin x < x$:
$\frac{\sin x}{x} < 1$؛ ومن
$x < \tan x = \frac{\sin x}{\cos x}$، بالضرب في
$\frac{\cos x}{x} > 0$: $\cos x < \frac{\sin x}{x}$. ولما
$x \to 0^+$، يكون $\cos x \to 1$: فبالحصر،
$\frac{\sin x}{x} \to 1$؛ و $\frac{\sin x}{x}$ \omterm{def:g11:func:parity}{زوجية}، إذن
النهاية من الجهتين هي $1$.

\textbf{4.} $\frac{1 - \cos x}{x^2} =
\frac{1 - \cos^2 x}{x^2 (1 + \cos x)} =
\left(\frac{\sin x}{x}\right)^2 \frac{1}{1 + \cos x}
\to 1 \times \frac12$.

\textbf{5.} $\sin'(0) = \lim_{h \to 0} \frac{\sin h - 0}{h} =
1$: فاشتقاق \omterm{def:g11:trigo:cossin}{الجيب} \omterm{def:g12:trigo:cossin}{وجيب التمام} كله ينبع من
هذه النهاية الواحدة. ولو قيست الزوايا بالدرجات، لكانت النهاية
$\frac{\pi}{180}$، ولجرّت كل \omterm{def:g12:deriv:derivative}{مشتقة} ذلك
الثابت: فالراديانات هي بالضبط الوحدة التي تجعل تحليل
التذبذب نظيفًا.

\textbf{6.} $y' = -R\omega\sin(\omega t - \varphi)$، و
$y'' = -R\omega^2\cos(\omega t - \varphi) = -\omega^2 y$.

\textbf{7.} $\omega = 2$. و $y(0) = A = 3$؛
و $y'(0) = 2B = 8$: إذن $B = 4$: $y = 3\cos 2t + 4\sin 2t$.

\textbf{8.} $R = \sqrt{9 + 16} = 5$؛ والدور
$\frac{2\pi}{2} = \pi$؛ والطور
$\varphi = \arctan\frac43 \approx 0.927$:
$y = 5\cos(2t - 0.927)$.

\textbf{9.} $5\cos(2t - \varphi) = 0$ أول مرة عندما يكون
$2t - \varphi = \frac\pi2$:
$t = \frac{\pi/2 + 0.927}{2} \approx 1.25$ s.

\textbf{10.} مع $\sin\theta \approx \theta$:
$\theta'' = -\frac gL \theta$: أي توافقية بالمقدار
$\omega = \sqrt{\frac gL}$، ودورها
$T = \frac{2\pi}{\omega} = 2\pi\sqrt{\frac Lg}$. ومن أجل
$L = 1$: $T \approx 2.006$ s. والدور لا يتعلق
بالسعة (من أجل التأرجحات الصغيرة) --- وهو تساوي الأزمنة عند غاليليو،
وهي الواقعة التي تجعل ساعات النواس ممكنة.

\textbf{11.} $L = g\left(\frac{T}{2\pi}\right)^2 =
\frac{9.81}{\pi^2} \approx 0.994$ m: فنواس الثانية
أقصر من متر خط الزوال بنحو $6$ mm. وقد اختارت الأكاديمية
خط الزوال (لأن الثقالة تتغير من مكان إلى مكان، فتخون
النواس)؛ ولو كان $g$ أكبر بشعرة، لكان مترنا
يدقّ.

\textbf{12.} $\cos\left(x + \frac\pi2\right) =
\cos x \cos\frac\pi2 - \sin x \sin\frac\pi2 = -\sin x$؛
$\sin\left(x + \frac\pi2\right) = \cos x$. إذن اشتقاق
\omterm{def:g11:trigo:cossin}{الجيب} يعطي \omterm{def:g11:trigo:cossin}{الجيب} منزاحًا ربع دورة (أي $\cos$)، ومرة أخرى
($-\sin$)، ومرة أخرى، فنعود إلى البيت في أربع خطوات:
فالاشتقاق يدير الدائرة.

\textbf{13.} بجمع نشرَي $\cos(a \mp b)$:
$\cos(a-b) + \cos(a+b) = 2\cos a\cos b$. ومع $a = b = x$:
$1 + \cos 2x = 2\cos^2 x$، أي
$\cos^2 x = \frac{1 + \cos 2x}{2}$.

\textbf{14.} مع $p = \frac{p+q}{2} + \frac{p-q}{2}$ و
$q$ نظيرتها، \omterm{def:g10:algebra:expand}{بالنشر} والجمع:
$\sin p + \sin q = 2\sin\frac{p+q}{2}\cos\frac{p-q}{2}$. ومن أجل
$440$ و $444$ Hz: نغمة عند $442$ Hz تُعدَّل شدتها
بالمقدار $\cos(2\pi \cdot 2t)$ --- فتبلغ الشدة ذروتها $4$
مرات في الثانية (وهي \omterm{def:g10:numbers:abs}{القيمة المطلقة} للغلاف): أي أربع ضربات
في الثانية، تتلاشى حين تتفق الشوكتان.

\textbf{15.} $\cos 3x = \cos 2x\cos x - \sin 2x \sin x =
(2\cos^2 x - 1)\cos x - 2\sin^2 x\cos x = 4\cos^3 x -
3\cos x$. وعند $x = 20^\circ$: $\cos 60^\circ = \frac12$، إذن
يحقق $c = \cos 20^\circ$ \omterm{def:g10:algebra:equation}{المعادلة} $8c^3 - 6c - 1 = 0$ --- وهي تكعيبية
لا حل لها بالجذور التربيعية: فإنشاء $20^\circ$، أي
تثليث $60^\circ$، خارج قدرة المسطرة والفرجار.

\textbf{16.} $R = \sqrt{9 + 27} = 6$؛
$\tan\varphi = \frac{3\sqrt3}{3} = \sqrt3$:
$\varphi = \frac\pi3$: $I(t) = 6\cos\left(100\pi t -
\frac\pi3\right)$: فالتيار الأعظمي $6$ أمبير، وتأخر الطور $60^\circ$.

\textbf{17.} $\sin^2 t = \frac{1 - \cos 2t}{2}$: فالدور
$\pi$. و $\sin 2t$ دورها $\pi$، و $\sin 3t$ دورها
$\frac{2\pi}{3}$: فيتكرر المجموع بعد أصغر مضاعف
مشترك، أي $2\pi$.

\textbf{18.} تذبذب دوره $1$ داخل الغلاف المنكمش
$\pm\eu^{-t/10}$: فكل تأرجحة أدنى بنحو $10\,\%$. وعند
$t = 10$ يكون الغلاف $\eu^{-1} \approx 0.368$ --- أي ثابت
\cref{pb:g12:exp:1} مرة أخرى. والمعادلات ذات الحلول
المخمَّدة كهذه ($y'' + a y' + b y = 0$) تُحل في فصل
المعادلات التفاضلية.

\textbf{19.} قوة تسحب راجعةً بالتناسب مع
الإزاحة تعطي $y'' = -\omega^2 y$؛ ويبيّن
\cref{exo:g12:trigo:9} أن حلول تلك \omterm{def:g10:algebra:equation}{المعادلة}
هي \emph{بالضبط} تراكيب $\cos$ و $\sin$ ---
فالوجود بالعرض، والوحدانية بالتمرين: فلا خيار للاهتزاز
إلا أن يكون جيبيًا.

\textbf{20.} نهاية واحدة ($\frac{\sin x}{x} \to 1$، بحصر
المساحات)، ومشتقتان ($\sin' = \cos$ و
$\cos' = -\sin$)، \omterm{def:g10:algebra:equation}{ومعادلة} واحدة ($y'' = -\omega^2 y$)، وعالم
من الساعات والأوتار والتيارات والمدّ. وأما النجمة: فمن أجل
التأرجحات الواسعة، $T = 4\sqrt{\frac Lg} \cdot
\frac{\pi/2}{M(1, \cos(\theta_0/2))}$ ---
\omterm{def:g11:stat:mean}{والمتوسط الحسابي} الهندسي في \cref{pb:g12:seq:1} يحسب
التكامل الإهليلجي في حفنة تكرارات: فمذكّرة غاوس
ونواس غاليليو على بعد صيغة واحدة.
\end{solution}
